\subsection{Order Delay}

OD is encoded to create an entirely new phase of attacks/deploys as already discussed in \ref{OD and phases}. For how it works read that. It is widely used and accepted in almost every vanilla SR and WR template. It's strength in these settings comes from the ability to simply out-delay an opponent on a very crucial turn, which often should result in better kill rates for you.
\\
\\
Besides the very straight forward usage it can have, it's often part of the timi-blockade sequence.
\\
\\
My advice would be to generally try and find good moves without using an OD card. In a Strat MME game where player A has already used their OD and player B hasn't, B always has an advantage. In contrast, you should not hoard the card over a turn where it would make a big difference. So in the end, there's a risk assessment factor of when, which conditions must be met, how, and the answer is it depends really. If you don't need to use it, then not use, it is strong enough to single handily win or lose you the game. For example my \href{https://www.warzone.com/MultiPlayer?GameID=38360784}{turn 6 usage} is simply not worth it, as I achieve nothing and live in fear over every other turn that I always might be outdelayed.